\documentclass[10pt]{article}

\usepackage[utf8]{inputenc}

\usepackage[T1]{fontenc}

\usepackage{amsmath}

\usepackage{amsfonts}

\usepackage{amssymb}

\usepackage[version=4]{mhchem}

\usepackage{stmaryrd}

\usepackage{bbold}



\begin{document}

$|a|=1$ существовало число $\delta=\delta(z, a)>0$ такое, что при $0<r<\delta$ выполняется неравенство



\[

u(z) \leqslant \frac{1}{2 \pi} \int_{0}^{2 \pi} u\left(z+r e^{i \varphi} a\right) d \varphi

\]



Для функций $u(z)$ класса $C^{2}(D)$ более удобен следуюший критерий: $u(z)$ есть П.ф. в $D$ тогда и только тогда, когда эрмитова форма (гессиан функции $\boldsymbol{u}$ )



\[

H((z ; u) a, \bar{a})=\sum_{j, k=1}^{n} \frac{\partial^{2} u}{\partial z_{j} \partial \bar{z}_{k}} a_{j} a_{k}

\]



неотрицательно определена в каждой точке $z \in D$.



Jum.: [1] В лад и м и р о в В. С., Методы теории функций многих комплексных переменных, М., $1964 . \quad$ Е.Д. Соломенцев. ПЛЮРИСУПЕРТАРМОНИЧЕСКАЯ ФУНКЦИЯ - сМ. Плюрисубгармоническая функиия.



ПОВЕРХНОСТЕЙ ГЕОМЕТРИЯ - раздел геометрии, изУчающий топологические и метрические свойства и внешнее строение (форму) поверхностей в евклидовых и, в общем случае, в римановых и псевдоримановых пространствах; для большей наглядности и краткости изложения ниже предполагается, что объемлющее пространство является евклидовым пространством $\mathbb{R}^{n}, n \geqslant 3$.



Само определение объекта исследования - поверхности - может быть различным. Наиболее часто встречаются следующие два варианта.



\begin{enumerate}

  \item Поверхность понимается как множество точек в $\mathbb{R}^{n}$, координаты к-рых удовлетворяют нек-рой системе уравнений

\end{enumerate}



\[

f_{i}\left(x^{1}, \ldots, x^{n}\right)=0,1 \leqslant i \leqslant m<n

\]



в таком виде поверхности встречаются, напр., как поверхности уровня, первые интегралы систем дифференциальных уравнений, поверхности раздела сред, области изменения переменных в задачах условного экстремума и т. д.



\begin{enumerate}

  \setcounter{enumi}{1}

  \item Поверхность можно понимать как отображение $f$ нек-рого многообразия $M$ в $\mathbb{R}^{n}$ (параметрически заданные поверхности); в этом случае поверхность представляется как параметризованное множество точек в $\mathbb{R}^{n}$, для K-рых так наз. внутренние координаты $u^{i}$ определяют положение точки $x$ на $M$ (в локальной карте), а через $x$ и $f(x)$ - положение точки на $S=f(M) \subset \mathbb{R}^{n}$.

\end{enumerate}



Среди геометрич. свойств поверхности различают внутренние и внешние свойства. Внутренняя геометрия поверхности $S$ описывает те ее свойства, к-рые зависят только от метрики поверхности. В свою очередь, метрики на $S$ имеют два возможных источника своего происхождения: а) поверхность $S$ задана заранее, напр., одним из указанных выше способов и метрика на ней тогда индуцируется через измерение длин кривых по закону метрики объемлющего пространства (расстояние между точками на $S$ определяется как inf длин кривых на $S$, соединяющих эти точки); б) поверхность $S$ получена изометрич. отображением риманова многообразия $M$ в $\mathbb{R}^{n}$; в этом случае говорят, что поверхность $S$ реализует в $\mathbb{R}^{n}$ метрику риманова многообразия $M$.



Проблема изометрич. погружения в $\mathbb{R}^{n}$ данного риманова многообразия является одной из центральных задач дифференциальной геометрии, и для классич. случая (погружение двумерного риманова многообразия в трехмерное евклидово пространство) решение ее дано с достаточной полнотой для метрик с кривизной $K>0, K<0$ и $K=0$, а для случаев знакопеременной кривизны известны лишь отдельные результаты. По поводу гладкости погружений есть общий результат, утверждающий, что поверхность $S$ (любой размерности) гладкости $C^{m, \alpha}, m \geqslant 2,0<\alpha<1$, в произвольном римановом пространстве той же гладкости $C^{m, \alpha}$ имеет метрику риманова многообразия гладкости $C^{m, \alpha}$, следовательно, данное риманово многообразие $M$ гладкости $C^{m, \alpha}$ нельзя изометрически погрузить в $\mathbb{R}^{n}$ поверхностью гладкости больше, чем $C^{m, \alpha}$.



Внутренняя геометрия поверхности не изменяется при ее изометрич. преобразованиях, в частности при ее изгибаниях. K ней относятся такие важные понятия, как длина кривой, угол между кривыми, площадь и объем, параллельный перенос векторов, геодезические и кратчайшие линии, градиент скалярного поля, тензор кривизны и т.д. Все эти величины и объекты можно определить, зная лишь первую квадратичную форму поверхности $S$, задаваемую во внутренних координатах $u^{i}$ в виде



\[

d s^{2}=g_{i j}(u) d u^{i} d u^{j}

\]



где дважды ковариантное тензорное поле $g_{i j}$ или вычисляется по уравнению $S$ в $\mathbb{R}^{n}$ или оно задано заранее, если $S$ определена как изометрич. погружение риманова многообразия $M$ в $\mathbb{R}^{n}$.



Внешняя геометрия поверхности связана с формой (внешним видом) расположения поверхности в пространстве. Важнейшей ее характеристикой является вторая квадратичная форма поверхности, выражающая, кратко говоря, разность ковариантного дифференцирования векторных полей на поверхности по закону метрики окружающего пространства и метрики самой поверхности. Эта форма определяет характер искривленности поверхности в пространстве, в частности, через еe коэффициенты можно вычислять кривизны линий на поверхности, ввести такие понятия, как главные кривизны, средняя кривизна, асимптотич. линии, эллипс кривизны и т. Д. Она позволяет также исследовать расположение поверхности по отношению к касательной плоскости и аппроқсимировать поверхность ее соприкасающимся параболоидом, в соответствии с классификацией точек поверхности как точки эллиптические (точки выпуклости), гиперболические (или седловые), параболические с обобщением этих типов точек на многомерные случаи.



Первая и вторая квадратичные формы определяют поверхность однозначно с точностью до движения в пространстве. Задавать, однако, эти формы заранее нельзя: их коэффициенты связаны так наз. основными уравнениями теории поверхностей.



Дополнительные требования на внешнее строение поверхности приводят к различного рода конкретным классам поверхностей - минимальные поверхности, выпуклые поверхности, развертывающиеся поверхности, поверхности вращения, линейчатые поверхности и т. Д. На существование и свойства поверхностей в целом существенное влияние имеет их топологич. строение. Топология поверхности важна также для теории функций и дифференциальных уравнений на поверхности.



Лит.: [1] Дубровин Б. А., Новиков С. П., Фоменко А. Т., Современная геометрия. Методы и приложения, М., 1986.



ПOBEРХНОСтНОГО НАТЯХЕнИЯ кOЭФФИИ. Сабитов. см. Минимальная поверхность.



ПОВОPОTA TОЧКА - точка $x_{0}$ для уравнения



\[

\varepsilon^{2} y^{\prime \prime}+q(x) y=0

\]



такая, что $q\left(x_{0}\right)=0$. П.т. называется простой, если $q^{\prime}\left(x_{0}\right) \neq 0$, и кратности $m$, если $q^{(j)}\left(x_{0}\right)=0, j \leqslant m-1$, $q^{(m)}\left(x_{0}\right) \neq 0$. Пусть на отрезке $I=[a, b]$ имеется ровно одна П. т. $x_{0}, a<x_{0}<b, q \in C^{\infty}(I)$ и она вещественна. Тогда асимптотика решений уравнения (1) при $\varepsilon \rightarrow+0$ (равномерно по $x \in I)$ выражается через функции Эйри при $m=1$ и через функции Бесселя при $m \geqslant 2$ (см. [1], [2]). Для уравнения



\[

\varepsilon^{2} y^{\prime \prime}+q(x, \varepsilon) y=0

\]



имеются два определения П. т.:



\begin{enumerate}

  \item $q\left(x_{0}, 0\right)=0$



  \item $q\left(x_{0}(\varepsilon), \varepsilon\right)=0 \quad$ и П.т. называется простой, если $q_{x}^{\prime}\left(x_{0}, 0\right) \neq 0$.



\end{enumerate}



Для простой П. т. асимптотика решений уравнений (2) в окрестности П.т. также выражается через функшии Эйри, для кратной (в общем случае) - через решения уравнения вида:



\[

\omega^{\prime \prime}+Q(x) \omega=0 \text {, }

\]





\end{document}